\section{A K3 surface and its stability wall}\label{sec:k3}

\subsection{An explicit surface}

In $\PP^4$ with coordinates $[x_0:\cdots:x_4]$, put
\begin{align}
q&=x_0x_3-x_1x_2+x_4^2,\nonumber\\
c_0&=x_0^3+x_0x_1^2+2x_0x_2^2+5x_1x_2x_3
       +4x_0x_1x_2+3x_1^2x_3+2x_2^2x_3,\label{eq:k3-polynomials}\\
c&=c_0+x_4(x_0^2+x_3^2+x_4^2),\qquad X=V(q,c).\nonumber
\end{align}

\begin{proposition}\label{prop:k3}
The surface $X$ is a smooth K3 surface. Its hyperplane section $x_4=0$
is the union of two smooth rational curves satisfying
\eqref{eq:intersection-intro}.
\end{proposition}

\begin{proof}
The quadratic form is nondegenerate, and $c$ is not divisible by $q$.
Smoothness can be checked by the following exact finite computation.
Over $\mathbb F_{101}$, take the ideal generated by $q,c$ and the ten
$2\times2$ minors of their Jacobian. On each affine chart $x_i=1$,
its reduced Gr\"obner basis is
\[
\begin{array}{c|ccccc}
i&0&1&2&3&4\\ \hline
\text{basis}&[1]&[1]&[1]&[1]&[1].
\end{array}
\]
All the polynomial inputs are printed in~\eqref{eq:k3-polynomials};
the accompanying calculation constructs the minors and uses exact
arithmetic in $\mathbb F_{101}$. Thus the geometric singular subscheme
in that fibre is empty, not merely its set of rational points.
The projective Jacobian subscheme over $\operatorname{Spec}\mathbb Z$
is proper. If its generic fibre were nonempty, its closed image would
contain the generic point and hence every prime, a contradiction at
101. The characteristic-zero Jacobian criterion proves smoothness.
Adjunction gives $K_X=\cO_X$, and the complete-intersection cohomology
sequence gives $H^1(X,\cO_X)=0$. Hence $X$ is K3.

On $q=x_4=0$ use Segre coordinates
$[x_0:x_1:x_2:x_3]=[au:av:bu:bv]$. The cubic restricts to $fg$, where
\begin{equation}\label{eq:two-curves}
f=a(u^2+v^2)+buv,\qquad
g=(a^2+2b^2)u+3abv.
\end{equation}
The two coefficients in each linear expression have no common zero.
Thus $D_1=V(f)$ and $D_2=V(g)$ are smooth rational graphs of bidegrees
$(1,2)$ and $(2,1)$ on $\PP^1\times\PP^1$.
On $a=u=1$, their intersection has
\[
-z^4+2z^2+2=0,\qquad b=-\frac{1+z^2}{z},\qquad z=v.
\]
There are four distinct nonzero roots. The fifth intersection has
$[a:b]=[u:v]=[0:1]$. In the chart $b=v=1$, substitution
$a=-u/(1+u^2)$ into $g$ has leading term $-u$, proving transversality
there. These exhaust the intersection number.

At a finite intersection, the normal derivative of $c$ is
$x_0^2+x_3^2=1+(1+z^2)^2=4(1+z^2)$ modulo the quartic, hence is
nonzero. At the fifth point it is one. The intersections are therefore
transverse also as curves on $X$. Their sum is the hyperplane section,
which is ample. Adjunction gives $D_i^2=-2$, and the five crossings give
$D_1D_2=5$. In particular $H^2=6$ and $HD_i=3$.
\end{proof}

There is also a noncomputational existence argument for the same
configuration. On the smooth quadric threefold $V(q)$, lift a product
section of bidegree $(3,3)$ from the hyperplane quadric. Vary its lift by
$x_4G$, where $G$ is quadratic. This linear variation is base-point-free
off the hyperplane, so a general member is smooth there by Bertini.
Along either component away from its transverse crossings, the
tangential derivative is nonzero. At the finitely many crossings, avoid
the finitely many affine hyperplanes on which the normal derivative
vanishes. The explicit equations above select one smooth member; no
assumption on the full Picard group of that member is needed.

\subsection{Bundles, obstruction and polarization}

For the Mukai convention $v(F)=(r,c_1,r+c_1^2/2-c_2)$,
\begin{equation}\label{eq:mukai-data}
v(L_i)=(1,D_i,0),\quad
v(E_0)=(3,2D_1+D_2,0),\quad c_2(E_0)=8,\quad v(E_0)^2=10.
\end{equation}
The class $2D_1+D_2$ pairs to one with $D_1$, so it is primitive.
The stable K3 moduli dimension is $v^2+2=12$
\cite{Muk84}. Nonemptiness in the chambers will follow from the
actual local construction, rather than from the dimension formula.

The line bundles have $H$-degree three and Euler characteristic one.
They are stable for every polarization, and their reduced $H$-Hilbert
polynomials agree. The differences $\pm(D_1-D_2)$ are nonzero classes
of $H$-degree zero and cannot be effective. Therefore
\[
\Hom(L_1,L_2)=\Hom(L_2,L_1)=0.
\]
Serre duality gives vanishing of both cross $\Ext^2$ groups.
Since $(D_1-D_2)^2=-14$, Riemann--Roch gives
\begin{equation}\label{eq:ext-dimensions}
\dim\Ext^1(L_1,L_2)=\dim\Ext^1(L_2,L_1)=5,\qquad
\Ext^1(L_i,L_i)=0.
\end{equation}
Serre-dual bases identify $\Ext^1(E_0,E_0)$ with $\cV$.
The quadratic Yoneda obstruction has blocks $-BA$ and $\tr(BA)$,
up to a common nonzero scalar; its zero locus is $BA=0$.
Section~\ref{sec:coordinates} establishes both harmonic normalization
of these bases and the convergent comparison with the full reduced
integrability germ.

The effective automorphism group is
$(GL_2\times GL_1)/\CC^*\simeq GL_2$.
The determinant-one group is a finite cover with an ineffective
$\mu_3$ center. A trivial effective stabilizer will mean simplicity of
the bundle, not literal absence of determinant-one scalar automorphisms.

For the ample real class $H_t=H+tD$, one has
\begin{align}
H_t^2&=6-14t^2,\nonumber\\
\mu_t(L_1)&=3-7t,\qquad \mu_t(L_2)=3+7t,\qquad
\mu_t(E_0)=3-\tfrac73t,\label{eq:slopes}\\
\mu_t(L_1)-\mu_t(E_0)&=-\tfrac{14}{3}t,\qquad
\mu_t(L_2)-\mu_t(E_0)=\tfrac{28}{3}t.\nonumber
\end{align}
Thus the slope ordering changes transversely. If $m$ is sufficiently
large, $mH\pm D$ is integral ample and
$c_1(E_0)(mH\pm D)=9m\mp7$ is coprime to three.
Equality of slopes with a subsheaf of rank one or two would force that
rank to be divisible by three. The standard implications between
slope and Gieseker stability then identify all the stable and
semistable loci for these rank-three invariants.
This arithmetic does not itself give a wall contraction; the local
HYM construction below supplies that identification.
